% -------------------------------------------------------------------------------------------------
%
%  Skeleton for semester and master thesis reports at the Laboratory for Software Technology.
%  Based on the skeleton provided by the Institute of Robotics and Intelligent Systems.
%
% -------------------------------------------------------------------------------------------------
%
% USAGE:      compile with PDFLaTeX
%
% HISTORY:    - written by Sascha A. Stoeter <stoeter@iris.ethz.ch>, www.stoeter.com, 02.06.2004
%             - modified by Martin Probst, 18.08.2004
%             - extended and adapted for use at LST by Oliver Trachsel, 2007-08-21
% -------------------------------------------------------------------------------------------------
\documentclass[11pt,a4paper]{book}
\usepackage{lstreport}

% -------------------------------------------------------------------------------------------------
% Add needed packages. Some generally useful packages are listed for
% your convenience.
% -------------------------------------------------------------------------------------------------
\usepackage{subfigure}                          % enable the use of subfigures
\usepackage[thickspace,thinqspace]{SIunits}     %
\usepackage[plainpages=false,pdfpagelabels]{hyperref}    % enable hyperlinks in pdf/ps Docs
\usepackage{listings}                           % to embed source code


% UTF8
\usepackage[utf8x]{inputenc} 

% -------------------------------------------------------------------------------------------------
% Select type of thesis
% -------------------------------------------------------------------------------------------------
\renewcommand{\thesistype}{Bachelor}
%\renewcommand{\thesistype}{Diploma}
%\renewcommand{\thesistype}{Master}

% -------------------------------------------------------------------------------------------------
% Set names
% -------------------------------------------------------------------------------------------------
\renewcommand{\thesisauthor}{Marcel Mohler}
\renewcommand{\thesisadvisor}{Zoltán Majó \\ Tobias Hartmann}

%OWN STUFF###############################################################
\usepackage[english]{babel}

% colors
%\definecolor{orange}{rgb}{1,0.5,0}

% Disable single lines at the start of a paragraph (Schusterjungen)

\clubpenalty = 10000

% Disable single lines at the end of a paragraph (Hurenkinder)

\widowpenalty = 10000
\displaywidowpenalty = 10000
 
% allows for colored, easy-to-find todos

%\newcommand{\todo}[1]{\textsf{\textbf{\textcolor{orange}{[[TODO: #1]]}}}}


%##########################################################################


% -------------------------------------------------------------------------------------------------
% Beginning of the main document body
% -------------------------------------------------------------------------------------------------
\begin{document}

% include all bibtex even if not cited
\nocite{*}

% -------------------------------------------------------------------------------------------------
% Front matter with title page, table of contents, and abstracts 
% -------------------------------------------------------------------------------------------------
\frontmatter

% Title page: set title and date 
\thesistitlepage{Profile Caching for the\\Java Virtual Machine
}{August 2015}

% Abstract must not be longer than one page per language. English and
% German abstracts are mandatory.
\chapter*{Introduction}
Virtual machines (VMs) like the Java Virtual Machine (JVM) are used as the execution environment of choice for many modern programming languages. 
VMs interpret a suitable intermediate language (e.g., Java Byte Code for the JVM) and provide the runtime system for application programs. VMs usually include a garbage collector, a thread scheduler, and interfaces to the host operating system. 
As interpretation of intermediate code is time-consuming, VMs usually include a \textit{Just-in-time} (JIT) compiler that translates frequently-executed functions or methods to \textit{native} machine code.
\\\\
The JIT compiler executes in parallel to a program's interpretation by the VM and, as a result, compilation speed is a critical issue in the design of a JIT compiler.
Unfortunately, it is difficult to design a compiler such that the compiler produces good or excellent code while limiting the resource demands of this compiler. The compiler requires storage, CPU cycles and even on a multi-core processor, compilation may slow down the execution of the application program.
\\\\
Consequently, most VMs adopt a multi-tier compilation system.
%The first tier usually is the interpretation of a method. If this method is executed frequently, the Tier-1 compiler translates this method into native code.
%The Tier-1 compiler implements only a small set of the know optimization techniques and as result, it had good compilation speed but the generated code is far from the output of an optimizing compiler. Such a compiler is usually the Tier-2 compiler, which takes a longer amount of time and produces optimized native code. To determine which methods should be compiled by the Tier-1 (or Tier-2) compiler, the VM profiles the execution of all application programs to identify “hot” methods.
At program startup, all methods are interpreted by the virtual machine (execution at Tier 0). The interpreter gathers execution statistics called \textit{profiles} and if a method is determined to be executed frequently, this method is then compiled by the Tier 1 compiler. Methods compiled to Tier 1 are then profiled further and based on these profiling information, some methods are eventually compiled at higher tiers.
One of the drawbacks of this setup is that for all programs, all methods start in Tier 0, with interpretation and profiling by the VM. However, for many programs the set of the most used methods does not change from one execution to another and there is no reason to gather profiling information again. 
\\\\
The main idea of this thesis is to cache these profiles from a prior execution to be used in further runs of the same program. 
Having these \textit{cached profiles} available avoids the JIT compiler to gather the same profiling information again. As well as allow the compiler to use more sophisticated profiles early in program execution and prevent recompilations when more information about the method is available. While this in general should not significantly influence the peak performance of the program, the hope is to decrease the time the JVM needs to achieve it, the so called \textit{warmup}.
\\\\
This thesis proposes a design and an implementation of a profile caching feature for \textit{HotSpot}, an open source Java virtual machine maintained and distributed by Oracle Corporation as well as a profound performance analysis using state-of-the-art benchmarks.


%\chapter*{Zusammenfassung}
%Kurzfassung der Arbeit.

% Table of contents
\tableofcontents

% -------------------------------------------------------------------------------------------------
% Main document body
% -------------------------------------------------------------------------------------------------
\mainmatter
\chapter{Introduction}
\label{s:Introduction}
Explain scope and structure of report.

\chapter{Great Work}
\label{s:GreatWork}

This and the following chapters detail the original work.

\chapter{Examples}
\label{s:Examples}

This chapter provides some additional hints and examples for the
layout and style of the thesis. It is worthwhile to look at the source
file \verb|Examples.tex| for this appendix to understand how it was
created.

\section{Tables}

Tables are left justified and the caption appears on top as seen in
Table~\ref{t:Translations}.

\begin{table}[ht]
\centering
\begin{tabular}{ll}
\hline
\textbf{English} & \textbf{German}\\
\hline
cell phone       & Handy\\
Diet Coke        & Coca Cola light\\
\hline
\end{tabular}
\caption[Translations]{\label{t:Translations}Translations.}
\end{table}

\section{Figures}

Figure~\ref{f:IRISlogo} shows a simple figure with a single picture
and Figure~\ref{f:SubfigureExample} shows a more complex figure
containing subfigures.

\begin{figure}[ht]
\centering
\includegraphics[width=.6\linewidth]{figures/LSTlogo}
\caption[LST logo]{\label{f:IRISlogo}LST logo.}
\end{figure}

\begin{figure}[ht]
\centering
\subfigure[ETH logo]{\includegraphics[height=12mm]{figures/ETHlogo}}\quad
\subfigure[LST logo]{\includegraphics[height=12mm]{figures/LSTlogo}}
\caption[Subfigure example]{\label{f:SubfigureExample}Two pictures as
  part of a single figure through the magic of the subfigure package.}
\end{figure}

\section{Units}

The SIUnits package provides nice spacing for units as demonstrated in
Table~\ref{t:SIUnits}. Use of the package also makes it easy to change
the style or even the unit text in the future.

\begin{table}[ht]
\centering
\begin{tabular}{ll}
\hline
\textbf{Output}   & \textbf{Command}\\
\hline
42m               & \verb|42m|\\
\unit{42}{\metre} & \verb|\unit{42}{\metre}|\\
42 m              & \verb|42 m|\\
\hline
\end{tabular}
\caption[Spacing for units]{\label{t:SIUnits}Spacing for units.}
\end{table}

\section{Source code}

The listings package provides tools to typeset source code
listings. It supports many programming languages and provides a lot of
formatting options.

\lstset{numbers=left, numberstyle=\tiny, stepnumber=1, numbersep=5pt}
\lstset{basicstyle=\ttfamily}
\lstset{frame=tb}

\begin{lstlisting}[float,caption=Example usage of the listing package,label=l:javaClass,language=Java]
class S {
   int f1 = 42;
   public S(int x) {
          f1 = x;
   }
}
\end{lstlisting}

Listing \ref{l:javaClass} shows an example listing. Code snippets can
also be inserted in normal text:
\verb$\lstinline|int f1 = 42;|$ gives \lstinline$int f1 = 42;$

\section{Miscellany}

\begin{description}

\item[Capitalization.] When referring to a named table (such as in the
  previous section), the word \emph{table} is capitalized. The same is
  true for figures, chapters and sections.

\item[Bibliography.] Use \verb|bibtex| to make your life easier and to
  produce consistently formatted entries.

\item[Contractions.] Avoid contractions. For instance, use ``do not''
  rather than ``don't.''

\item[Style guide.] A classic reference book on writing style is
  Strunk's \emph{The Elements of Style} \cite{Strunk-ElementsOfStyle}.

\end{description}

% -------------------------------------------------------------------------------------------------
% Appendices (if needed)
% -------------------------------------------------------------------------------------------------
\appendix
\chapter{Extra Stuff}
\label{s:ExtraStuff}

Additional material such as long mathematical derivations.

% -------------------------------------------------------------------------------------------------
% Bibliography
% -------------------------------------------------------------------------------------------------
\addcontentsline{toc}{chapter}{Bibliography}
\bibliography{bibliography}

\end{document}
